\section{The exact cost of blocking torsion specialization}
\label{sec:torsion}

\subsection{Definitions and the sharp theorem}
Fix a prime \(p\), and let \(v\) be the valuation on \(\Cp\), normalized by \(v(p)=1\), with \(v(0)=+\infty\). Put
\[
 e_h=(p-1)p^{h-1},\qquad r_h=e_h^{-1},\qquad
 Z_h=\{\zeta-1:\ord(\zeta)=p^h\}\quad(h\ge1).
\]
Every element of \(Z_h\) has valuation \(r_h\), and \(\#Z_h=e_h\).
For a nonzero polynomial \(g(T)=\sum_j a_jT^j\), define
\[
 G_u(g)=\min_j\{v(a_j)+ju\}\quad(u\ge0),\qquad
 \Delta_\rho(g)=G_\rho(g)-G_0(g).
\]
The quantity \(\Delta_\rho(g)\) is unchanged by multiplying \(g\) by a nonzero scalar. It measures how much the Gauss valuation changes between the coefficient norm and the specified positive radius. It is nonnegative.

Fix \(h_0\ge1\), \(\rho\ge r_{h_0}\), and write \(c=p^{1-h_0}\).
A polynomial \emph{blocks} the shells \(h_0,\ldots,H\) if
\[
 v(g(z))>G_\rho(g)
 \quad\text{for every }z\in Z_h,\quad h_0\le h\le H.
\]
The comparison radius \(\rho\) is fixed throughout. It is not replaced by the radius \(r_h\) of the point currently being tested.

\begin{theorem}[Exact obstruction budget]
\label{thm:obstruction}
For \(H\ge h_0\), the minimum of \(\Delta_\rho(g)\) among all nonzero polynomials in \(\Cp[T]\) blocking the shells \(h_0,\ldots,H\) is
\[
 \boxed{\ (H-h_0+1)p^{1-h_0}\ }.
\]
The minimum is attained by a monic integral polynomial with roots among those torsion points. If \(N=H-h_0+1\), one such polynomial has degree \(p^N-1\).
\end{theorem}

\begin{corollary}[Optimal universal endpoint]
\label{cor:endpoint}
Let \(B\ge0\) be finite. If \(\Delta_\rho(g)\le B\), then some \(h\) and \(z\in Z_h\) satisfy
\[
 v(g(z))\le G_\rho(g),\qquad
 h_0\le h\le H_*(B):=h_0+\floor{Bp^{h_0-1}}.
\]
The endpoint is best possible uniformly in degree and coefficient field. In particular, the same conclusion holds for integral \(g\) with \(G_\rho(g)\le B\).
\end{corollary}

The unrestricted coefficient field is part of the sharpness statement. A restriction to a fixed field, such as \(\Qp\), can impose additional Galois symmetry on the roots and improve the bound; such refinements were already studied in the preceding reports.

\subsection{The exact one-shell capacity}
The following estimate is the local combinatorial ingredient. Its floor-sum form occurs in the prior shell-minimax reports \cite{shellprior,samplingprior}; a complete proof is given here so that the new aggregate argument is self-contained.

\begin{lemma}[One-shell capacity]
\label{lem:shell}
Suppose \(m<e_h\) and \(v(\alpha_i)=r_h\) for \(1\le i\le m\), with repetitions allowed. Then some \(z\in Z_h\) satisfies
\[
 \sum_{i=1}^m\bigl(v(z-\alpha_i)-r_h\bigr)
 \le U_h(m):=\frac1{e_h}
 \sum_{j=1}^{h-1}e_j\floor{\frac{m}{e_j}}.
 \tag{S}
\]
The bound is sharp. With \(m=e_h\), placing one root at every point makes the minimum excess infinite.
\end{lemma}

\begin{proof}
Choose a primitive \(p^h\)-th root \(\xi\). For each \(\alpha_i\), select a unit residue \(b_i\bmod p^h\) maximizing
\(v(\xi^{b_i}-1-\alpha_i)\).
For every unit residue \(a\), ultrametricity gives
\[
 v(\xi^a-1-\alpha_i)\le v(\xi^a-\xi^{b_i}).
\]
If \(a=b_i\), the right side is infinite and the inequality is automatic. Otherwise, a strict reverse inequality would force
\(v(\xi^{b_i}-1-\alpha_i)<v(\xi^a-1-\alpha_i)\), contradicting maximality. Thus projecting each root to a nearest torsion point can only increase the potential at every endpoint.

View the units modulo \(p^h\) as the leaves of a tree. There are \(p-1\) first branches, and every subsequent branch has \(p\) children. Mark the leaves \(b_i\), with multiplicity. Choose a least-occupied first branch and then a least-occupied child at every stage. At depth \(j\), the chosen class contains at most
\[
 \floor{m/e_j}
\]
marks. This follows inductively from
\(\lfloor\lfloor m/e_j\rfloor/p\rfloor=\lfloor m/e_{j+1}\rfloor\).
Since \(m<e_h\), the final leaf \(a\) is unmarked.

For \(a\ne b_i\), put \(t_i=v_p(a-b_i)\), interpreted modulo \(p^h\); then \(0\le t_i<h\), and
\[
 v(\xi^a-\xi^{b_i})-r_h
 =\frac{p^{t_i}-1}{e_h}
 =\frac1{e_h}\sum_{j=1}^{t_i}e_j.
\]
Summing first over \(i\) and then over depths counts the marks in the selected classes. The preceding occupancy bounds prove (S).

For sharpness, mark the first \(m\) positive integers prime to \(p\). In each block of \(e_j\) such integers, every unit residue modulo \(p^j\) appears once. Hence every unit class has at least \(\lfloor m/e_j\rfloor\) marks. An unmarked endpoint has potential at least the right side of (S); a marked endpoint has infinite potential. The upper bound just proved supplies equality somewhere.
\end{proof}

The nearest-point projection is why the lemma covers arbitrary \(\Cp\)-roots, not just cyclotomic roots. Neither algebraicity of the coefficients nor a Galois-stable root multiset is assumed.

\subsection{Root accounting and the indispensable earlier-shell savings}
Factor
\(g=b\prod_\alpha(T-\alpha)\), with multiplicities. Gauss multiplicativity gives
\[
 G_u(g)=v(b)+\sum_\alpha\min(u,v(\alpha)).
 \tag{G}
\]
This includes roots of negative valuation and the root \(0\), whose valuation is infinite.

Let \(m_h\) be the number of roots of valuation exactly \(r_h\). Every root contributes nonnegatively to \(G_\rho-G_0\), while a root on a selected shell contributes exactly \(r_h\). Therefore
\[
 \Delta_\rho(g)\ge\sum_{h=h_0}^H\frac{m_h}{e_h}.
 \tag{B}
\]
For \(z\in Z_h\), roots off shell \(h\) have
\(v(z-\alpha)=\min(r_h,v(\alpha))\). Consequently
\[
 v(g(z))-G_\rho(g)
 =
 \sum_{v(\alpha)=r_h}\bigl(v(z-\alpha)-r_h\bigr)-s_h,
 \qquad
 s_h:=G_\rho(g)-G_{r_h}(g).
 \tag{V}
\]
The slack \(s_h\) contains the savings supplied by roots on earlier shells:
\[
 s_h\ge
 \sum_{k=h_0}^{h-1}(r_k-r_h)m_k.
 \tag{E}
\]
Roots outside the selected shells add nonnegative slack and cannot make blocking cheaper. Keeping the full quantity in (E), instead of discarding savings through a positive-part operation, is the step that changes the global minimax problem.

\subsection{Proof of the exact lower bound}
\begin{proof}[Proof of the lower bound in Theorem~\ref{thm:obstruction}]
We prove by induction that blocking through shell \(h\) requires
\[
 m_h\ge M_h:=(p-1)p^{h-h_0}.
 \tag{M}
\]
Put \(k=h-h_0\), so \(M_h=e_{k+1}\). Assuming the earlier lower bounds,
\[
 s_h\ge
 \sum_{i=h_0}^{h-1}(r_i-r_h)M_i
 =kc-\frac{p^k-1}{e_h}
 =\frac{kM_h-(p^k-1)}{e_h}.
 \tag{L}
\]
The last expression equals \(U_h(M_h-1)\). Indeed,
\begin{align*}
 e_hU_h(M_h-1)
 &=\sum_{j=1}^{k+1}(M_h-e_j)\\
 &=(k+1)M_h-(p^{k+1}-1)\\
 &=kM_h-(p^k-1).
\end{align*}
When \(k+1=h\), the extra \(j=h\) term is zero, so this also covers \(h_0=1\). When \(k=0\), both sides are zero, supplying the induction base.

If \(m_h<M_h\), then \(m_h<e_h\). Lemma~\ref{lem:shell}, monotonicity of \(U_h\), and (L) give an endpoint whose on-shell excess is at most \(s_h\). Equation (V) then gives \(v(g(z))\le G_\rho(g)\), contrary to blocking. Thus (M) holds.

Since \(M_h/e_h=c\), summing (B) yields
\[
 \Delta_\rho(g)\ge(H-h_0+1)c.
\]
\end{proof}

\subsection{A sharp construction}
\begin{proof}[Proof of attainment in Theorem~\ref{thm:obstruction}]
On each shell \(h_0\le h\le H\), choose a primitive root \(\xi_h\) and mark the first \(M_h\) positive unit exponents. Equivalently, mark the exponents
\[
 1\le a\le p^{h-h_0+1},\qquad p\nmid a.
\]
Let \(g\) be the monic product of all the factors
\(T-(\xi_h^a-1)\).
All roots have positive valuation, so the coefficients are integral and \(G_0(g)=0\). Formula (G) gives
\[
 G_\rho(g)=\sum_{h=h_0}^H M_hr_h=(H-h_0+1)c.
\]

At shell \(h\), the slack from earlier shells is exactly (L); later shells add no slack. If \(h_0=1\), then \(M_h=e_h\), and \(g\) vanishes at every selected endpoint. If \(h_0\ge2\), then \(M_h<e_h\), and the balanced marks give minimum on-shell excess
\[
 U_h(M_h)=(k+1)c.
\]
This is strictly greater than the slack
\(kc-(p^k-1)/e_h\). In fact their difference is
\[
 \frac{p^{k+1}-1}{e_h}>0.
\]
Equation (V) therefore proves strict blocking at every point.

Finally, with \(N=H-h_0+1\), the degree is
\[
 \sum_{k=0}^{N-1}(p-1)p^k=p^N-1.
\]
The roots, hence all coefficients, lie in a finite cyclotomic extension of \(\Qp\).
\end{proof}

\begin{proof}[Proof of Corollary~\ref{cor:endpoint}]
Take \(N=\lfloor B/c\rfloor+1\). Blocking those \(N\) shells would require
\(\Delta_\rho(g)\ge Nc>B\). Thus success occurs by the stated endpoint.

Conversely, if \(n=\lfloor B/c\rfloor\ge1\), the sharp construction on \(n\) shells has budget \(nc\le B\) and blocks through \(h_0+n-1\). The endpoint cannot be decreased. If \(n=0\), the first allowed shell is already the least possible endpoint. This also explains the strict threshold when \(B/c\) is an integer: \(n\) shells can still fail, while \(n+1\) force success.

For integral \(g\), \(G_0(g)\ge0\), so \(\Delta_\rho(g)\le G_\rho(g)\le B\).
\end{proof}

\subsection{Power series and simultaneous aggregate comparisons}
For a nonzero series
\(F=\sum_{\nu\in\mathbf N^s}a_\nu\mathbf T^\nu\)
whose coefficient valuations are bounded below, define
\[
 G_0(F)=\inf_\nu v(a_\nu),\qquad
 G_\rho(F)=\min_\nu\{v(a_\nu)+\rho|\nu|\},\qquad
 \Delta_\rho(F)=G_\rho(F)-G_0(F).
\]
The infimum in \(G_0\) need not be attained. The minimum in \(G_\rho\) is attained, since the coefficient valuations have a lower bound and \(\rho|\nu|\to\infty\).
The series converges whenever all coordinate valuations are positive.

\begin{theorem}[Bounded series in any finite dimension]
\label{thm:series}
If \(\Delta_\rho(F)\le B\), there are \(h_j\in[h_0,H_*(B)]\) and \(z_j\in Z_{h_j}\) such that
\[
 v(F(z_1,\ldots,z_s))\le G_\rho(F).
\]
The endpoint is independent of \(s\). The coordinates need not have equal orders.
\end{theorem}

\begin{proof}
First suppose \(F\) is a polynomial, and set \(A=G_\rho(F)\) and \(a=G_0(F)\). Write
\[
 F=\sum_\beta f_\beta(T_1)\mathbf T_{\mathrm{rest}}^\beta
\]
and choose \(\beta\) attaining \(A\), so
\(G_\rho(f_\beta)+\rho|\beta|=A\).
Every coefficient of \(f_\beta\) has valuation at least \(a\), hence
\[
 \Delta_\rho(f_\beta)\le G_\rho(f_\beta)-a\le A-a\le B.
\]
The scalar theorem gives \(z_1\) with
\(v(f_\beta(z_1))\le G_\rho(f_\beta)\).
The polynomial in the remaining variables has Gauss valuation at most \(A\), and every coefficient still has valuation at least \(a\), because \(v(z_1)>0\). Its variation is therefore at most \(B\). It is nonzero since the selected coefficient has finite valuation. Repeating the argument selects all coordinates without increasing the budget.

For a bounded series, keep the real number \(a=\inf_\nu v(a_\nu)\) as a lower bound; do not normalize by a scalar whose valuation might not exist. Put \(H=H_*(B)\), and choose \(D\) so that
\[
 a+(D+1)\min(\rho,r_H)>A.
\]
The truncation \(P\) to total degree at most \(D\) has
\(G_\rho(P)=A\) and \(G_0(P)\ge a\), so
\(\Delta_\rho(P)\le B\).
At every point of the finite torsion grid, the omitted tail has valuation strictly greater than \(A\). Apply the polynomial result to \(P\); at the chosen point \(v(P)\le A\), so strict ultrametricity gives \(v(F)=v(P)\le A\).
\end{proof}

\begin{corollary}[Finite packets]
\label{cor:packets}
Let \(F_1,\ldots,F_n\) be nonzero bounded series in the same finite set of variables. If
\(\sum_i\Delta_\rho(F_i)\le B\), there is a common torsion tuple in the same grid with
\[
 \sum_{i=1}^n v(F_i(\mathbf z))
 \le \sum_{i=1}^nG_\rho(F_i).
\]
In particular, for integral series it suffices that
\(\sum_iG_\rho(F_i)\le B\).
Separate witness points \(\mathbf x_i\) with coordinate valuations at least \(\rho\) and
\(\sum_i v(F_i(\mathbf x_i))\le B\)
also suffice in the integral case.
\end{corollary}

\begin{proof}
Apply Theorem~\ref{thm:series} to the product \(\prod_iF_i\).
Gauss valuations at positive radii are additive. Their limits as the radius tends to zero are \(G_0\), so \(G_0\) and \(\Delta_\rho\) are additive as well. The product is nonzero and bounded. The finite resulting valuation implies no evaluated factor vanishes. Finally,
\(G_\rho(F_i)\le v(F_i(\mathbf x_i))\) at each witness.
\end{proof}

\subsection{Why the aggregate qualification is necessary}
The following counterexample is inherited from the earlier exact report \cite{exactprior}, not a new claim. Let \(p=2\), fix \(h_0\), choose a primitive \(2^h\)-th root \(\xi_h\) for each \(h_0\le h\le L\), and put
\[
 f_h(T)=T-(\xi_h-1).
\]
For \(\rho\ge r_{h_0}\),
\[
 \sum_{h=h_0}^L G_\rho(f_h)=\sum_{h=h_0}^L r_h<2r_{h_0}.
\]
Yet every \(z=\zeta-1\in Z_h\) satisfies
\(v(f_h(z))>r_h=G_\rho(f_h)\): two primitive roots of the same \(2\)-power order have a quotient of strictly smaller order, unless they are equal, when the valuation is infinite.

Thus a bounded aggregate budget can defeat every \emph{individual} exact comparison on arbitrarily many shells. The product can still satisfy the exact aggregate comparison, because savings on earlier factors offset a loss on the currently active factor. The same example blocks comparison with the moving Gauss target \(G_{r_h}\). Neither stronger objective is asserted by Theorem~\ref{thm:series}.

\subsection{Extension to Lubin--Tate torsion}
\label{subsec:lt}
Let \(K/\Qp\) be finite, \(\pi\) a uniformizer, \(q\) its residue-field cardinality, and \(w\) the extended valuation normalized by \(w(\pi)=1\). Let \(\mathcal F\) be a one-dimensional Lubin--Tate formal \(\Ocal_K\)-module for \(\pi\). Write \(Z_h^{\mathcal F}\) for its points of exact \(\pi^h\)-torsion. Here the formal parameter itself is evaluated; there is no additional subtraction of \(1\).

\begin{corollary}[Exact Lubin--Tate selection]
\label{cor:lt}
Replace \(p\) by \(q\), \(v\) by \(w\), and \(Z_h\) by \(Z_h^{\mathcal F}\) in Theorems~\ref{thm:obstruction} and \ref{thm:series}. Their conclusions remain valid:
\[
 \min\Delta_\rho(g)=Nq^{1-h_0},\qquad
 H_*(B)=h_0+\floor{Bq^{h_0-1}},
\]
provided
\(\rho\ge((q-1)q^{h_0-1})^{-1}\).
The scalar sharp construction has degree \(q^N-1\).
\end{corollary}

\begin{proof}
The standard Lubin--Tate structure theorem gives
\(\mathcal F[\pi^h]\simeq\Ocal_K/\pi^h\), so the primitive layer has
\(e_h=(q-1)q^{h-1}\) points \cite{lt}.
The defining series satisfies
\([\pi](T)=\pi T+\cdots\) and
\([\pi](T)\equiv T^q\pmod\pi\).
Its Newton polygon shows that primitive \(\pi\)-torsion has valuation \(1/(q-1)\). Inductively, if \([\pi](x)\) has primitive level \(h-1\), then \(w(x)<1/(q-1)\), and the \(T^q\)-term uniquely gives the minimum valuation, hence
\(w([\pi](x))=q\,w(x)\).
Thus primitive level \(h\) has valuation
\(r_h=1/e_h\).

To compare ordinary and formal differences, write
\[
 \mathcal F(X,[-1]_{\mathcal F}Y)=(X-Y)U(X,Y).
\]
The series on the left vanishes at \(X=Y\); its linear term is \(X-Y\), so \(U\) has integral coefficients and constant term \(1\). At points of positive valuation, \(w(U)=0\).
For a generator \(t_h\) and distinct unit residues \(a,b\bmod\pi^h\),
\[
 w([a]t_h-[b]t_h)
 =w([a-b]t_h)
 =r_{h-w(a-b)}.
\]
The unit-residue tree therefore has \(q-1\) first branches and \(q\) subsequent branches, with exactly the distances used in Lemma~\ref{lem:shell}. Least-occupied branches give the same floor sums. Residue-balanced representatives modulo \(\pi^{h-h_0+1}\) attain them. Root accounting, the induction, and truncation then carry over unchanged.
\end{proof}

This is a quantitative selection statement. Serban's work already studies torsion rigidity for Lubin--Tate formal groups \cite{serban}; the corollary does not claim a new Manin--Mumford classification.

